\documentclass[11pt]{article}
\usepackage[a4paper,margin=27mm]{geometry}
\usepackage{amsmath,amssymb,amsthm,mathtools}
\usepackage[T1]{fontenc}
\usepackage{lmodern}
\usepackage{microtype}
\usepackage[colorlinks=true,linkcolor=blue,citecolor=blue,urlcolor=blue]{hyperref}
\newtheorem{theorem}{Theorem}
\newtheorem{lemma}[theorem]{Lemma}
\newtheorem{corollary}[theorem]{Corollary}
\theoremstyle{remark}
\newtheorem{remark}[theorem]{Remark}
\newcommand{\ddc}{\mathrm{d}\mathrm{d}^{\mathrm c}}
\newcommand{\dc}{\mathrm{d}^{\mathrm c}}
\newcommand{\C}{\mathbb C}
\newcommand{\R}{\mathbb R}
\newcommand{\Pp}{\mathbb P}
\newcommand{\FS}{\mathrm{FS}}
\newcommand{\MA}{\operatorname{MA}}
\newcommand{\negpart}[1]{[#1]_-}
\title{Sharp damping thresholds for rotationally invariant\newline
potentials on a Hopf threefold}
\author{Research manuscript}
\date{October 7, 2026}
\begin{document}
\maketitle
\begin{center}
\small Working draft. Mathematical review and independent reconstruction pending.
\end{center}
\begin{abstract}
On the standard Hopf threefold we determine the damping thresholds in the
full class of normalized potentials invariant under the unitary rotations
of the last two coordinates. Every individual weighted non-pluripolar
measure is finite exactly when $\sqrt t\,W(t)$ is bounded; this also gives
a uniform mass bound. Uniform tails, uniform capacity absolute continuity
and capacity-to-weak continuity instead require $\sqrt t\,W(t)\to0$.
An exact positive fiber flux and a pointwise semiconvex depth estimate prove
sufficiency for all bounded and unbounded potentials in this class. A
faster-scale heat construction gives the necessary conditions, a single
infinite measure, and an actual capacity-convergent sequence whose weighted
measures charge a pole divisor in the limit at the critical regime.
The unrestricted damping problem remains open.
\end{abstract}

\section{The statement and its scope}
Let
\[
 X=(\C^3\setminus\{0\})/\langle z\mapsto \mathrm e^{-1}z\rangle,
 \qquad \ddc=\frac{\mathrm i}{2\pi}\partial\bar\partial,
 \qquad \dc=\frac{\mathrm i}{4\pi}(\bar\partial-\partial).
\]
On the cover put $s=\log|z|^2$, $\alpha=\ddc s$,
$\beta=\mathrm ds\wedge\dc s$, and $\omega=\alpha+\beta$.
Then $\alpha=\pi^*\omega_{\FS}$ for projectivization
$\pi:X\to\Pp^2$, with $\int_{\Pp^2}\omega_{\FS}^2=1$ and
$\int_X\omega^3=6$. All potentials below are genuine quasi-plurisubharmonic
functions with their canonical upper representatives. For a smooth bounded
potential the extended non-pluripolar product is its ordinary smooth product.

\begin{theorem}\label{thm:core}
There are $A,B,D,\eta_0>0$ such that for every $0<\eta<\eta_0$ there exist
$u_\eta\in C^\infty(X,\R)$ and a measurable $K_\eta\subset X$ satisfying
\begin{gather*}
 \sup_Xu_\eta=0,\qquad \omega+\ddc u_\eta>0,
 \qquad 0\leq-u_\eta\leq B\eta^{-2}\quad\text{on }X,\\
 A\eta^{-2}\leq-u_\eta\leq B\eta^{-2}\quad\text{on }K_\eta,
 \qquad
 \int_{K_\eta}(\omega+\ddc u_\eta)^3\geq D\eta^{-1}.
\end{gather*}
\end{theorem}

\begin{corollary}\label{cor:weights}
Let $W:[0,\infty)\to(0,\infty)$ be nonincreasing. A uniform bound for the
weighted masses
\[
 \int_X W(-u)\MA^{\rm ext}_\omega(u)
\]
over all normalized $\omega$-psh potentials requires
\[
 \sup_{T\gg1}\sqrt T\,W(T)<\infty.
\]
Uniform deep tails over that class require $\sqrt T\,W(T)\to0$.
Consequently, for eventually decreasing weights
$W_{p,a}(T)=(1+T)^{-p}[\log(\mathrm e+T)]^{-a}$ tending to zero,
uniform mass bounds fail when $p<1/2$, or when $p=1/2$ and $a<0$.
Uniform deep tails also fail when $p=1/2$ and $a=0$.
\end{corollary}

\begin{theorem}\label{thm:individual}
Suppose $W$ is positive and nonincreasing, finite on compact intervals, and
\[
 \sup_{T\gg1}\sqrt T W(T)=\infty.
\]
Then there is a single genuine normalized $\omega$-psh potential $U$ such that
\[
 \int_XW(-U)\MA^{\rm ext}_\omega(U)=\infty.
\]
Thus universal individual finiteness fails in the same power/logarithmic
regimes where Corollary~\ref{cor:weights} gives unbounded uniform masses.
\end{theorem}

Changing a weight on a fixed compact interval does not affect these
obstructions. The heat construction is credited to Xia--Zhang; the faster
scale, correction and limiting arguments are proved below rather than
importing their final construction at a different parameter.

Let $\mathcal R$ denote the \emph{entire} class of normalized genuine
$\omega$-psh potentials invariant under
\[
 (z_0,z')\longmapsto(z_0,Az'),\qquad A\in U(2),\quad z'=(z_1,z_2).
\]
This is a globally defined holomorphic symmetry preserving $\omega$.
No reduction of arbitrary potentials to $\mathcal R$ is assumed. Write
$\mu_u=\MA^{\rm ext}_\omega(u)$ and
$\nu_{W,u}=W(-u)\mu_u$, with zero measure on the pole set.
The capacity in this paper is exactly
\[
 \operatorname{Cap}_\omega(B)=
 \sup_{-1\leq v\leq0,\ \omega+\ddc v\geq0}
       \int_B(\omega+\ddc v)^3.
\]
For a positive $C^1$ nonincreasing $W$ tending to zero put
\[
 Q(W)=\sup_{t\geq0}(1+\sqrt{2t})W(t),\qquad
 Q_T(W)=\sup_{t\geq T}(1+\sqrt{2t})W(t).
\]

\begin{theorem}\label{thm:sharp}
Within $\mathcal R$ the following assertions hold.
\begin{enumerate}
\item Every individual $\nu_{W,u}$ is finite if and only if
$\sup_{t\gg1}\sqrt t\,W(t)<\infty$. This is also equivalent to a uniform
total-mass bound. In this case
\[
 \sup_{u\in\mathcal R}\nu_{W,u}(X)\leq12Q(W).
\]
If the condition fails there is one $u\in\mathcal R$ with infinite measure.
\item Uniform deep tails hold if and only if $\sqrt t\,W(t)\to0$.
For all $u\in\mathcal R$ and $T\geq0$,
\[
 \nu_{W,u}(\{-u>T\})\leq12Q_T(W).
\]
The inequality is understood in the extended sense when its right side
is infinite.
\item Uniform absolute continuity relative to $\operatorname{Cap}_\omega$
holds if and only if $\sqrt t\,W(t)\to0$.
\item The assignment $u\mapsto\nu_{W,u}$ is a map into finite measures
and is continuous from capacity convergence in $\mathcal R$ to weak
measure convergence if and only if $\sqrt t\,W(t)\to0$.
When $\sqrt t\,W(t)$ is bounded but does not tend to zero, there are
actual smooth $u_j\in\mathcal R$ converging in capacity to an actual
$g\in\mathcal R$, whose weighted measures have a weak cluster measure
charging the pole divisor of $g$.
\end{enumerate}
For eventually decreasing power/logarithmic weights, modified smoothly
on a compact interval if necessary, finiteness and uniform total bounds
hold exactly for $p>1/2$ or $p=1/2,a\geq0$. The other three properties
hold exactly for $p>1/2$ or $p=1/2,a>0$.
\end{theorem}

The unrestricted sufficient condition
$\int_0^\infty(1+t)(-W'(t))\,\mathrm dt<\infty$ remains known background.
Theorem~\ref{thm:sharp} improves it on the declared invariant class only.
The general optimal-rate question and the interval between the necessary
square-root scale and existing unrestricted sufficient scales remain open.

\section{Heat estimates and the local positivity test}
On the affine chart $U_0\simeq\C^2$ write
$\xi_j=z_j/z_0$ and $w=\log z_0$ modulo
$\Lambda=\mathbb Z+2\pi\mathrm i\mathbb Z$.
The fibers are $E=\C/\Lambda$, with probability area
$\mathrm dA=\mathrm i\,\mathrm dw\wedge\mathrm d\bar w/(4\pi)$.
For $\xi\ne0$ set
\[
 x=\log|\xi|^2,\qquad \mu=\frac{\mathrm e^x}{1+\mathrm e^x},
 \qquad \mu_x=\mu(1-\mu).
\]
For a smooth real $V(x,w)$ independent of $x$ at the negative end and
constant in both variables at the positive end, smooth global extension
holds. The positivity test of \cite[Lemma 2.2]{XZ} is
\begin{equation}\label{eq:test}
 1+V_{w\bar w}>0,\qquad \mu+V_x>0,\qquad
 Q[V]:=V_{xx}+\mu-
          \frac{|\mu+V_{x\bar w}|^2}{1+V_{w\bar w}}>0.
\end{equation}
At $\xi=0$ only the first condition remains. Indeed, the matrix of
$2\pi(\omega+\ddc V)$ in the horizontal tangential, radial, and vertical
frame used there is
\[
 \begin{pmatrix}
 (\mu+V_x)/\mu&0&0\\
 0&(\mu+V_{xx})/[\mu(1-\mu)]&
                      (\mu+V_{x\bar w})/\sqrt{\mu(1-\mu)}\\
 0&(\mu+V_{wx})/\sqrt{\mu(1-\mu)}&1+V_{w\bar w}
 \end{pmatrix}.
\]
The Schur complement gives \eqref{eq:test}. Our $\omega$ and $\ddc$ are
both the source's divided by $2\pi$; scalar potentials and positivity are
therefore unchanged.

Define the fiber heat kernel and potential by
\begin{equation}\label{eq:heat}
 k_t(w)=\frac2t\sum_{\lambda\in\Lambda}\mathrm e^{-|w-\lambda|^2/t},
 \qquad
 F_\eta(t,w)=-\eta\int_t^\infty(k_r(w)-1)\,\mathrm dr.
\end{equation}
Unfolding gives $\int_Ek_t\,\mathrm dA=1$.
Its Fourier coefficients are
$\exp(-\pi^2|p+\mathrm iq/(2\pi)|^2t)$, so
\begin{gather}
 (F_\eta)_t=(F_\eta)_{w\bar w}=\eta(k_t-1),\qquad
 (\log k_t)_{w\bar w}\geq-1/t,\label{eq:heatfacts}\\
 \|\partial_t^aF_\eta(t,\cdot)\|_{C^m(E)}
       \leq C_{a,m}\eta\mathrm e^{-t/4}\quad(t\geq1).
 \label{eq:largetime}
\end{gather}
The logarithmic Hessian equals $-1/t$ plus the variance of $w-\lambda$,
under the Gaussian lattice weights, divided by $t^2$.
The smallest nonzero Fourier eigenvalue is $1/4$, proving the last bound.

\begin{lemma}\label{lem:values}
Uniformly in $\eta$ and $w$,
$F_\eta(t,w)\leq C\eta$ for $t>0$, and
$F_\eta(t,w)\geq-2\eta\log(1/t)-C\eta$ for $0<t\leq1$.
Uniformly on $\sqrt t\leq|w|\leq2\sqrt t$ near the origin of $E$, as $t\to0$,
\begin{align*}
 F_\eta(t,w)&=2\eta\log t+O(\eta),\\
 t k_t(w)&=2\mathrm e^{-|w|^2/t}+O(\mathrm e^{-1/(Ct)}),\\
 t\partial_{\bar w}\log k_t(w)&=-w+O(\mathrm e^{-1/(Ct)}).
\end{align*}
\end{lemma}
\begin{proof}
For the upper bound use $k_r\geq0$ on $(0,1]$ and \eqref{eq:largetime}.
The nonnegative Fourier coefficients imply $k_r(w)\leq k_r(0)$, and Gaussian
summation gives
$k_r(0)\leq2/r+C r^{-1}\mathrm e^{-1/(2r)}$ for $r\leq1$.
Its integrable error proves the lower bound.
On the indicated annulus the nonzero-lattice terms and their derivatives
are exponentially small relative to the zero-lattice term, proving the last
two identities. For the first, split the defining integral at a fixed small
$r_0$. On $[t,r_0]$ the nonzero-lattice contributions have bounded integral,
and the difference between $(2/r)\mathrm e^{-|w|^2/r}$ and $2/r$ has integral
at most $2|w|^2/t\leq8$. The remaining integral is uniformly bounded.
\end{proof}

We also use \cite[Lemmas 3.2--3.3]{XZ}. For fixed
$0<\kappa<1/(200C_0)$, $0<\eta<1/10$, $0<t\leq1$, $1/2\leq\mu\leq1$, and
\[
 \sigma^2=\frac{2\kappa}{\eta}\log\frac1{\eta t},\quad
 h=\eta k_t,\quad
 \zeta=\frac\sigma\mu t\partial_{\bar w}\log k_t,
\]
their estimate is
\begin{equation}\label{eq:square}
 \kappa tk_t-
 \frac{\mu^2(1-\eta)h}{1-\eta+h}
       |\zeta-(1-\eta)^{-1}|^2
 \leq C\kappa\exp\left(-\frac1{C\sigma^2t}\right).
\end{equation}
It follows from
$tk_t\leq C_0\exp(-|t\partial_{\bar w}\log k_t|^2/(C_0t))$.
To recall its application, for a nearest lattice point with distance $d$,
Gaussian summation bounds $tk_t$ by $C\mathrm e^{-d^2/t}$ and
$|t\partial_{\bar w}\log k_t|$ by $C(d+\mathrm e^{-1/(Ct)})$.
This proves the latter estimate. If $\sigma^2t\geq1/100$, \eqref{eq:square}
is immediate after enlarging $C$. Otherwise the square absorbs $\kappa tk_t$
when $|\zeta-(1-\eta)^{-1}|\geq1/4$, using
$\kappa t(1+h)<\eta/100$. In the remaining case $|\zeta|\geq3/4$ and the
Gaussian estimate applies. Constants are uniform for smaller positive
$\kappa$, which will be used below.

\section{The faster scale and global positivity}
Fix $\kappa$ as above, also with $\kappa<1/100$. Choose a small fixed $c>0$
later and set
\begin{gather*}
 \varepsilon=\mathrm e^{-c\eta^{-3}},\qquad
 q_0=\sqrt{c\eta^{-3}+\log(1/\eta)},\qquad
 a=\sqrt{\kappa/(2\eta)},\\
 q(x)=q_0-\int_0^{ax}\chi(s)\,\mathrm ds,\qquad
 t(x)=\eta^{-1}\mathrm e^{-q(x)^2},\\
 b=(q_0-1/2)/a,\qquad x_*=(q_0+1/2)/a,\qquad
 Z_\eta(x,w)=\vartheta(q(x))F_\eta(t(x),w).
\end{gather*}
Here $\chi$ is a fixed smooth nondecreasing cutoff from zero to one on
$[0,1]$, with $\int_0^1\chi=1/2$; $\vartheta$ is another fixed smooth cutoff
from zero to one on $[0,1]$. Then $Z_\eta=F_\eta(\varepsilon,\cdot)$ for
$x\leq0$, $Z_\eta=0$ for $x\geq x_*$, and
\[
 \eta b,\eta x_*\longrightarrow\ell:=\sqrt{2c/\kappa},\qquad
 1+(Z_\eta)_{w\bar w}>1-\eta.
\]
The source uses a different initial heat time. We recheck the positivity
estimates at this new scale rather than import its final proposition.

On $[0,b]$ we have $\vartheta=1$.
With $\sigma=(\log t)_x$, direct differentiation gives the source's exact
square identity
\begin{equation}\label{eq:identity}
\begin{split}
 Q[Z_\eta]={}&\sigma^2t\{h[1+t(\log k_t)_{w\bar w}]-\eta\}
       +\sigma_xtF_t+\mu(1-\mu)-\frac{\mu^2\eta}{1-\eta}\\
 &+\frac{\mu^2(1-\eta)h}{1-\eta+h}
             |\zeta-(1-\eta)^{-1}|^2.
\end{split}
\end{equation}
The contribution involving $h[1+t(\log k_t)_{w\bar w}]$ is nonnegative.
On $[0,a^{-1}]$,
\[
 \sigma=2aq\chi(ax),\quad
 \sigma_x=-\frac\kappa\eta\chi(ax)^2+\frac\kappa\eta q\chi'(ax),
 \quad t\leq\varepsilon\mathrm e^{q_0}.
\]
This upper bound for $t$ is smaller than every power of $\eta$.
Thus $\kappa q\chi't\leq\eta$ and $\sigma^2t\leq\eta^2$.
Use \eqref{eq:square} with $\kappa\chi^2$ where $\chi>0$, and $F_t\geq-\eta$
for the $\chi'$ term. It follows that $\negpart{Q[Z_\eta]}\leq C\eta$.
Where $\chi=0$ also $\chi'=0$, and
$Q[Z_\eta]=\mu-\mu^2/(1-\eta+h)>0$, since $x\leq a^{-1}=o(1)$.

On $[a^{-1},b]$, $\sigma=2aq$, $\sigma_x=-\kappa/\eta$,
$\mathrm dx=-\mathrm dq/a$. For $t\leq1$, \eqref{eq:identity} and
\eqref{eq:square} give
\[
 \negpart{Q[Z_\eta]}\leq
 2\eta+2\kappa q^2t+C\kappa\exp\left(-\frac\eta{C\kappa q^2t}\right).
\]
For $t\geq1$, \eqref{eq:largetime} gives instead
$\negpart{Q[Z_\eta]}\leq C\kappa q^2t^2\mathrm e^{-t/4}$.
Put $L=\log(1/\eta)$. Integrating the $q^2t$ term uses
$\int_{\sqrt L}^\infty q^2\mathrm e^{-q^2}\,\mathrm dq\leq C\sqrt L\eta$.
For the exponential term split at $2\sqrt L$: the initial interval has
length $O(\sqrt L)$, and on the rest
$\eta^2\mathrm e^{q^2}/q^2\geq\mathrm e^{q^2/4}$ for small $\eta$.
Its remaining integral is uniformly bounded even with upper endpoint infinity.
For large time use
$\mathrm dx=\sqrt{\eta/(2\kappa)}\,\mathrm dt/(tq)$,
$q\leq\sqrt L$, and the integrability of $t\mathrm e^{-t/4}$.
Including the constant term and initial interval yields
\begin{equation}\label{eq:defect}
 \int_0^b\max_{w\in E}\negpart{Q[Z_\eta](x,w)}\,\mathrm dx
       \leq C\eta b+C\sqrt{\eta L}\leq C\ell+o(1).
\end{equation}
All constants here are independent of $q_0,c,\eta$ with fixed $\kappa$.
The bound is small after choosing $c$, although it need not tend to zero.

Also $\sup\negpart{(Z_\eta)_x}=O(\sqrt{\eta L})$:
for $t\leq1$ use $F_t\geq-\eta$ and
$t\sqrt{L+\log(1/t)}\leq\sqrt L$; for $t\geq1$ use \eqref{eq:largetime}.
On $[b,x_*]$, $0\leq q\leq1$ and $(\mathrm e\eta)^{-1}\leq t\leq\eta^{-1}$,
so differentiated cutoffs and \eqref{eq:largetime} give
\[
 \sup_w\bigl(|(Z_\eta)_x|+|Q[Z_\eta]-Q[0]|\bigr)
       \leq C\mathrm e^{-1/(5\mathrm e\eta)}.
\]
Choose $c$ so that $\ell<1/(10\mathrm e)$. Since
$Q[0]=\mu(1-\mu)\geq\mathrm e^{-x_*}/4$ here, this implies $Q[Z_\eta]>0$
on the cutoff interval. For $x\leq0$ use
$Q[Z_\eta]\geq\mu(1-\mu/(1-\eta))>0$, and for $x\geq x_*$ use $Z_\eta=0$.
On $[-2,-1]$ the same formula gives $Q[Z_\eta]\geq2/27$.

The continuous defect $d_\eta(x)=\max_w\negpart{Q[Z_\eta]}$ is therefore
supported in $[0,b]$. Choose a smooth nonnegative $r_\eta$, supported in
$(-1/2,b+1/2)$, with $r_\eta>d_\eta$ on $[0,b]$ and
$m_\eta:=\int r_\eta\leq\int d_\eta+\eta$.
Uniform approximation and a positive bump give this choice for each fixed
$\eta$; no uniform derivative bound is needed.
Fix a smooth probability bump $\psi$ supported in $(-2,-1)$ and put
\[
 g_\eta(x)=\int_{-\infty}^x(x-s)(r_\eta(s)-m_\eta\psi(s))\,\mathrm ds,
 \qquad V_\eta=Z_\eta+g_\eta.
\]
The support order implies $-m_\eta\leq g_\eta'\leq0$;
$g_\eta=0$ at the negative end and is constant at the positive end.
Furthermore $g_\eta\leq0$ and $|g_\eta|\leq m_\eta(b+3)=O_c(\eta^{-1})$.
Shrink the fixed $c$ so that \eqref{eq:defect} ensures
\[
 m_\eta<\min\{\mu(-2)/2,\,1/8,\,1/(27\|\psi\|_\infty)\}
\]
for all sufficiently small $\eta$.
Now $Q[V_\eta]=Q[Z_\eta]+r_\eta-m_\eta\psi>0$ everywhere:
use the majorant on $[0,b]$, the margin $2/27$ on $(-2,-1)$, and the
established positivity elsewhere.
The vertical coefficient is above $1-\eta$. For the horizontal coefficient,
derivatives vanish for $x\leq-2$; on $[-2,0]$ use $\mu\geq\mu(-2)>m_\eta$;
on $[0,\infty)$ use $\mu\geq1/2$ and $(Z_\eta)_x\geq-o(1)$.
Thus \eqref{eq:test} gives a smooth globally strictly positive potential.

On $x<-3$, $V_\eta=F_\eta(\varepsilon,\cdot)$ has fiber mean zero, so
$0\leq\sup_XV_\eta\leq C\eta$ by Lemma~\ref{lem:values} and $g_\eta\leq0$.
Set $u_\eta=V_\eta-\sup_XV_\eta$.
Where the cutoff is nonzero $t\geq\varepsilon$, giving
\begin{equation}\label{eq:globaldepth}
 0\leq-u_\eta\leq2c\eta^{-2}+O_c(\eta^{-1})+O(\eta).
\end{equation}
Multiplication by a cutoff between zero and one preserves this lower bound
for $Z_\eta$, including where $F_\eta$ is positive.

\section{Mass at the asserted depth}
Let $K_\eta$ consist of all radial angular directions and
\[
 q_0/2\leq q(x)\leq3q_0/4,\qquad
 \sqrt{t(x)}\leq|w|\leq2\sqrt{t(x)},
\]
using the small lift near zero in $E$.
This lies in the linear transport region, with $\vartheta=1$,
$x\asymp_c\eta^{-1}\to\infty$, $\sigma=2aq\asymp_c\eta^{-2}$, and
superexponentially small $t$. Lemma~\ref{lem:values} gives, uniformly,
\[
 0<c_1\leq tk_t\leq3,\quad h\to\infty,\quad
 \zeta\to0,\quad \eta\sigma^2t\to0,\quad \mu\to1.
\]
In \eqref{eq:identity} the first $h$ contribution is nonnegative,
$\sigma_xtF_t=-\kappa t(k_t-1)\geq-3\kappa$, and the completed square
tends to one. Its other negative terms tend to zero.
Thus $Q[Z_\eta]\geq1/4$ for small $\eta$. On the core $\psi=0$ and
$g_\eta''=r_\eta\geq0$, so $Q[V_\eta]\geq1/4$. Also
\[
 \mu+(V_\eta)_x\geq c_2\eta\sigma,\qquad
 1+(V_\eta)_{w\bar w}\geq c_1\eta/t.
\]
Indeed $(Z_\eta)_x=\eta\sigma(tk_t-t)$ and $\mu+g_\eta'>0$ here;
the vertical inequality follows from $1-\eta+\eta k_t$.

Let $\mathrm d\Omega$ denote probability angular area in $\C^2$.
The radial Fubini--Study measure has ball mass $\mu^2$, hence
$\omega_{\FS}^2=2\mu^2(1-\mu)\,\mathrm dx\,\mathrm d\Omega$.
Since $\omega^3=6\omega_{\FS}^2\mathrm dA$, the determinant of the displayed
positivity matrix gives the exact density
\begin{equation}\label{eq:density}
 (\omega+\ddc V)^3
 =12(\mu+V_x)(1+V_{w\bar w})Q[V]\,
           \mathrm dx\,\mathrm d\Omega\,\mathrm dA.
\end{equation}
The fiber annulus has area $3t/2$.
Integrating the preceding lower bounds and using
$\sigma\,\mathrm dx=-2q\,\mathrm dq$ yields
\[
 \int_{K_\eta}(\omega+\ddc u_\eta)^3
 \geq c_3\eta^2\int_{q_0/2}^{3q_0/4}2q\,\mathrm dq
 =c_4\eta^2q_0^2\geq c_5\eta^{-1}.
\]
This places actual mass on the core; no inference from total energy alone
or limiting unbounded product is used.
The annular value estimate also gives
\[
 -u_\eta=2\eta q^2+O_c(\eta^{-1})+O(\eta\log(1/\eta)).
\]
As $\eta q_0^2=c\eta^{-2}+\eta\log(1/\eta)$, this is bounded above and below
by fixed positive multiples of $\eta^{-2}$ on the core.
Together with \eqref{eq:globaldepth} this proves Theorem~\ref{thm:core}.
All these potentials are bounded: truncations beyond their individual depth
equal them everywhere, so the defining extended product equals
\eqref{eq:density} exactly.

\begin{proof}[Proof of Corollary~\ref{cor:weights}]
On $K_\eta$, monotonicity of $W$ gives
\[
 \int_XW(-u_\eta)\MA^{\rm ext}_\omega(u_\eta)
       \geq D\eta^{-1}W(B\eta^{-2}).
\]
Every sufficiently large $T$ equals $B\eta^{-2}$ for some sufficiently small
$\eta$, proving the necessary bound.
Also $K_\eta\subset\{-u_\eta>A\eta^{-2}/2\}$, so uniform deep tails force
this lower bound to tend to zero. The power and logarithmic conclusions
follow from $\sqrt T W_{p,a}(T)\sim T^{1/2-p}(\log T)^{-a}$.
\end{proof}

\section{A single infinite weighted measure}
\begin{proof}[Proof of Theorem~\ref{thm:individual}]
First select the correction with additional local control. Put
\[
 x_{\rm pre}=(q_0+1/2-2\sqrt L)/a.
\]
Then $t(x_{\rm pre})=\eta^3$, $x_{\rm pre}\to\infty$, and
$b-x_{\rm pre}=O(\sqrt{\eta L})\to0$.
For a fixed large $K$ the preceding defect estimates give
$d_\eta\leq(K/2)\eta$ on $[0,x_{\rm pre}]$:
on the linear portion $q^2t\leq4L\eta^3$, and the exponential error is
$O(\exp(-1/(C\kappa\eta^2L)))=o(\eta)$.
On the initial portion use the previously proved $C\eta$ bound.
The integral of $d_\eta$ on $[x_{\rm pre},b]$ tends to zero by the same
integrated estimates and this interval's length.

Choose a fixed left cutoff, zero for $x\leq-1/2$ and one for $x\geq-1/4$,
and a translated right cutoff, one for $x\leq x_{\rm pre}$ and zero for
$x\geq x_{\rm pre}+1/4$.
Let $r^0_\eta$ be $K\eta$ times their product. It strictly exceeds $d_\eta$
on $[0,x_{\rm pre}]$. The continuous function
$e_\eta=\max(d_\eta-r^0_\eta,0)$ is supported in $[x_{\rm pre},b]$.
Add a smooth nonnegative majorant $s_\eta>e_\eta$ on that interval,
supported in $(x_{\rm pre}-1/4,b+1/2)$, with
$\int s_\eta\leq\int e_\eta+\eta$.
Set $r_\eta=r^0_\eta+s_\eta$. Then
\[
 r_\eta>d_\eta\quad\text{on }[0,b],\qquad
 m_\eta=\int r_\eta\leq K\eta b+O(\sqrt{\eta L})+O(\eta).
\]
After shrinking the fixed $c$, this is bounded by a fixed $m_0$ within all
the positivity tolerances above. The construction and smooth-core theorem
therefore hold for this choice, with fixed $A,B,D$.
On every fixed compact interval $I\Subset(0,\infty)$, for sufficiently small
$\eta$ we have $r_\eta=K\eta$ and exactly
\begin{equation}\label{eq:localcorrection}
 g_\eta(x)=-m_\eta(x-\mu_\psi)
              +K\eta(x^2/2+b_0x+b_1),\qquad
 \mu_\psi=\int s\psi(s)\,\mathrm ds.
\end{equation}
Here $b_0,b_1$ depend only on the fixed left cutoff.
Globally $0\leq-g_\eta(x)\leq m_0(x+2)$ for $x\geq0$.

Let $G(w)=-\int_0^\infty(k_r(w)-1)\,\mathrm dr$ for $w\ne0$ on $E$.
Gaussian integration gives the uniform bound
\begin{equation}\label{eq:offpolevalue}
 -F_\eta(t,w)\leq
 C\eta\bigl(1+\log^+(1/\operatorname{dist}_E(w,0)^2)\bigr),\quad t>0.
\end{equation}
Indeed the nearest Gaussian has integral bounded by
$C+2\log^+(1/\operatorname{dist}_E(w,0)^2)$ on $(0,1]$;
the other lattice terms have bounded integral and large times decay.
The same bound applies after the cutoff.
On a fixed $I\Subset(0,\infty)$, $\vartheta=1$ for small $\eta$ and
$t_\eta(x)\leq\exp(-c/(2\eta^3))$.
Every fixed $x$ derivative of $t_\eta$ has this bound times a polynomial
in $\eta^{-1}$. On compact fiber sets away from zero,
\[
 F_\eta(t_\eta(x),w)/\eta\longrightarrow G(w)
                   \quad\text{in }C^\ell\text{ for every fixed }\ell.
\]
To verify derivatives, subtract $G$ from $G_t=F_\eta(t,\cdot)/\eta$:
$G_t-G=\int_0^t(k_r-1)\,\mathrm dr=-t+O(\exp(-d^2/(Ct)))$ when
$d$ stays positive. Differentiated Gaussian errors have only polynomial
factors in $t^{-1}$, and repeated $x$ differentiation preserves their decay.
Together with \eqref{eq:localcorrection} and $0\leq\sup V_\eta\leq C\eta$,
this gives uniform $C^\ell$ bounds for $u_\eta$ on each such fixed compact
product set, for all sufficiently small $\eta$. The threshold may depend
on the set and derivative order.

Fix $a_j=2^{-j-2}$, $j\geq1$, so $\sum a_j=1/4$.
Choose scales $\eta_j\downarrow0$ recursively and write
$u_j=u_{\eta_j}$, $K_j=K_{\eta_j}$.
Require the minimum $x$ on $K_j$ to exceed every earlier $x_{*,i}$ by two.
Choose a relatively compact product neighborhood $O_j$ of $K_j$, with
positive $x$ interval and fiber closure away from zero, on which the earlier
$u_i$ are constant. The $K_j$ are pairwise disjoint.
With $E_j=\sum_{i<j}a_i\|u_i\|_\infty$, impose
\begin{gather*}
 E_j\leq a_j\eta_j^{-2},\qquad \eta_j\leq a_j,\\
 a_j\eta_j\leq2^{-j}a_i\eta_i\quad(i<j).
\end{gather*}
For each earlier $O_i$, require the preceding uniform derivative bounds to
hold for $u_j$ through order $j$. All are finite constraints at stage $j$,
achievable by making $\eta_j$ smaller. They ensure that the future series
$\sum_{k>i}a_ku_k$ converges in $C^\infty$ on compact subsets of $O_i$:
for every fixed derivative order the sufficiently late terms have a uniform
bound there, and $\sum a_k<\infty$.

These constraints also give $\sum_{k>j}a_k\eta_k\leq a_j\eta_j$.
On $K_j$ the coordinate and distance estimates are
\[
 x\leq C_c\eta_j^{-1},\qquad
 \log^+(1/\operatorname{dist}_E(w,0)^2)\leq C_c\eta_j^{-3}.
\]
Using \eqref{eq:offpolevalue}, the global slope bound on $g_k$, and the
supremum constants, the future depth is at most
\[
 C_c\eta_j^{-3}\sum_{k>j}a_k\eta_k
       +C_c\eta_j^{-1}\sum_{k>j}a_k
 \leq C_c a_j\eta_j^{-2}+C_c\eta_j^{-1}/4.
\]
The last term is at most a fixed multiple of $a_j\eta_j^{-2}$ since
$\eta_j\leq a_j$. Including the earlier and current terms, there is a fixed
$C_*>0$, independent of the scales, such that
\begin{equation}\label{eq:sumdepth}
 -\sum_i a_iu_i\leq C_*a_j\eta_j^{-2}\quad\text{on }K_j.
\end{equation}
Fix this constant before the scale choice. Since $\sqrt T W(T)$ is unbounded,
at each stage choose arbitrarily large $T_j$ such that
\[
 \frac D{\sqrt{C_*}}a_j^{5/2}\sqrt{T_j}W(T_j)\geq1,
 \qquad \eta_j=\sqrt{C_*a_j/T_j},
\]
and so small a scale that all finite preceding constraints hold.
Such values occur arbitrarily far out because $W$ is finite on compact
intervals. Thus the recursive choice is possible.

The smooth nonpositive partial sums $v_N=\sum_{j\leq N}a_ju_j$ decrease, and
\[
 \omega+\ddc v_N
 =(1-\sum_{j\leq N}a_j)\omega
       +\sum_{j\leq N}a_j(\omega+\ddc u_j)\geq\tfrac34\omega.
\]
Normalized $\omega$-psh potentials have a uniform $L^1$ bound on compact
connected $X$: add fixed local quadratic potentials, apply submean at a
maximum point in one inner coordinate ball, and propagate the bound along
a finite connected ball cover using the positive volumes of its overlaps.
This argument uses no mass bound. Hence $\int(-v_N)\leq C\sum a_j<\infty$.
The decreasing limit $v$ is $L^1$ finite and not identically $-\infty$.
Locally, after adding the same quadratic potential, the decreasing limit is
genuine psh; it retains its canonical upper representative.
Distributional positivity passes to the limit with margin $3\omega/4$.
Its supremum is finite and nonpositive.

On $O_j$ the earlier terms are constant and the future series converges
smoothly. Thus $v$ is smooth there and
$\omega+\ddc v\geq a_j(\omega+\ddc u_j)$ pointwise.
The determinant inequality for positive Hermitian matrices gives on $K_j$
\[
 \MA^{\rm ext}_\omega(v)\geq a_j^3(\omega+\ddc u_j)^3.
\]
This is an actual local smooth product: on a compact neighborhood of $K_j$
all sufficiently deep canonical truncations equal $v$.
Set $U=v-\sup_Xv$. Since $\sup v\leq0$, $-U\leq-v$, and its local currents
are unchanged. Monotonicity of $W$ and \eqref{eq:sumdepth} give
\[
 \int_{K_j}W(-U)\MA^{\rm ext}_\omega(U)
 \geq D a_j^3\eta_j^{-1}W(C_*a_j\eta_j^{-2})
 =\frac D{\sqrt{C_*}}a_j^{5/2}\sqrt{T_j}W(T_j)\geq1.
\]
The disjoint $K_j$ prove infinite total mass for this single potential.
They lie in its smooth locus, so no mass is assigned to poles.
\end{proof}

\section{Positive fiber flux and the weighted upper estimate}
We now work with an arbitrary smooth $V\in\mathcal R$, with $V\leq0$.
In the affine chart it has the form $V(x,w)$. The matrix calculation above
applies at every finite $x$, without assuming the special endpoint behavior
of the heat family. Put
\[
 A_V=\mu+V_x,\quad B_V=1+V_{w\bar w},\quad
 R_V=\mu+V_{xx},\quad C_V=\mu+V_{x\bar w},\quad
 \Delta_V=B_VR_V-|C_V|^2.
\]
Semipositivity gives $A_V,B_V,R_V,\Delta_V\geq0$ and, without dividing
by $B_V$,
\begin{equation}\label{eq:general-density}
 \MA_\omega(V)=12 A_V\Delta_V\,
       \mathrm dx\,\mathrm d\Omega\,\mathrm dA.
\end{equation}
The omitted divisor charts have zero mass for a smooth top form.

\begin{lemma}\label{lem:flux}
For every smooth invariant genuine $\omega$-psh potential,
\[
 \int_{\R\times E}\Delta_V\,\mathrm dx\,\mathrm dA=1.
\]
If in addition $V\leq0$, then
\[
 |V_x|^2\leq-2V,\qquad 0\leq A_V\leq1+\sqrt{-2V}.
\]
\end{lemma}
\begin{proof}
Set $m(x)=\int_EV\,\mathrm dA$ and
$e(x)=\int_E|V_w|^2\,\mathrm dA$.
Periodic integration by parts gives
\[
 \begin{split}
 \int_E V_{w\bar w}V_{xx}\,\mathrm dA
    &=-\operatorname{Re}\int_EV_{\bar w}V_{xxw}\,\mathrm dA,\\
 e''&=2\int_E|V_{xw}|^2\,\mathrm dA
       +2\operatorname{Re}\int_EV_{\bar w}V_{xxw}\,\mathrm dA.
 \end{split}
\]
The mean of every fiber derivative vanishes. Expanding $\Delta_V$ yields
\[
 \int_E\Delta_V\,\mathrm dA
   =m''+\mu(1-\mu)-\tfrac12e''.
\]
Thus $f(x)=\log(1+\mathrm e^x)+m(x)-e(x)/2$ is convex.

Both $m$ and $e$ are smooth global functions on the projective base:
local fiber coordinates change by holomorphic translations, preserving
fiber Haar area, the vertical derivative and the Dirichlet integral.
They are $U(2)$-invariant. Smoothness and radial invariance at $\xi=0$
give $m',e'\to0$ at the negative end. At the positive end use a normal
coordinate $z_0/z_1$ to the hyperplane $H=\{z_0=0\}$ on the base.
The invariant base functions are constant along this hyperplane and
smooth across it; their $x$ derivatives tend to zero. This does not
require that the pointwise potential be constant on its elliptic fibers
at that end. Consequently $f'(-\infty)=0$ and $f'(+\infty)=1$.
Integrating the nonnegative $f''$ proves the flux identity.

For each fixed $w$, $R_V\geq0$ implies $V_{xx}\geq-1$ on all of $\R$.
Hence for every real $h$,
\[
 V(x+h,w)\geq V(x,w)+hV_x(x,w)-h^2/2.
\]
Taking $h=V_x(x,w)$ and using $V\leq0$ proves the derivative bound.
Finally $\mu\leq1$ proves the bound for $A_V$.
\end{proof}

Combining the positive density with Lemma~\ref{lem:flux} proves
\begin{equation}\label{eq:smooth-upper}
 \int_XW(-V)\MA_\omega(V)\leq12Q(W),\qquad
 \int_{\{-V>T\}}W(-V)\MA_\omega(V)\leq12Q_T(W).
\end{equation}
The same argument works for any nonnegative continuous weight with finite
$Q$, without its monotonicity. There is no bound on unweighted mass in
this argument; the positive flux supplies the finite reference integral.

\section{Bounded approximation in the original capacity}
We give the passage to all genuine bounded potentials and the local
product argument explicitly. Only bounded local Bedford--Taylor products,
integration by parts, quasi-continuity and polar-set capacity zero are
used from \cite{BT}; no global unbounded Hopf product theorem is invoked.

\begin{lemma}\label{lem:capacity}
The original capacity is locally comparable, on a finite cover by inner
coordinate balls, with relative bounded-psh Bedford--Taylor capacities on
larger balls. Every bounded genuine $u\in\mathcal R$ with $-L\leq u\leq0$
has smooth genuine invariant approximants $u_\epsilon$ with the same
bounds converging in $\operatorname{Cap}_\omega$. If uniformly bounded
genuine potentials converge in that capacity, their bounded products,
and those products multiplied by continuous functions of their depths,
converge weakly.
\end{lemma}
\begin{proof}
For the upper capacity comparison add a fixed positive quadratic $\rho$
with $\ddc\rho\geq\omega$ to every global test $v\in[-1,0]$.
The closed positive cube of $\ddc(v+\rho)$ dominates
$(\omega+\ddc v)^3$. Normalize its bounded local potential range to
$[0,1]$ and use the relative capacity definition.
For the reverse comparison take a local psh test $w\in[-1,0]$ on the
larger ball of radius $R$, and an inner radius $r<R$.
Choose $a>0$ small so $a\ddc|z|^2\leq\omega$, then
$0<\epsilon<a(R^2-r^2)$ small. Select $-1<C_0<0$ with
\[
 -aR^2<C_0<-\epsilon-ar^2.
\]
The potential $\max(\epsilon w-a|z|^2,C_0)$ equals the constant near
the outer boundary and the first branch on the inner ball. Extend by
$C_0$ globally. It is a genuine global test in $[-1,0]$ and its current
on the inner ball dominates $\epsilon\ddc w$. This proves the reverse
comparison with the fixed factor $\epsilon^{-3}$. These comparisons
transfer local quasi-continuity and polar-set capacity zero to the
original capacity.

The compact group $G=U(3)\times(\R/\mathbb Z)$ acts transitively and
holomorphically on $X$ by unitary transformations and positive scalar
multiplication. It preserves $\omega$. With a nonnegative smooth central
approximate identity $\rho_\epsilon$ on $G$, put
\[
 u_\epsilon(x)=\int_Gu(gx)\rho_\epsilon(g)\,\mathrm dg.
\]
Transitive compact-group convolution is smooth, averaging preserves the
current inequality, and centrality makes convolution commute with the
$U(2)$ action. Thus $u_\epsilon\in\mathcal R$ except that its supremum
need only be nonpositive; \eqref{eq:smooth-upper} requires only this.
The original capacity is $G$-invariant by transporting its genuine
bounded tests and local products.

Choose an open $O$ of arbitrarily small capacity such that $u$ is
continuous on compact $X\setminus O$. For $g$ close to the identity,
$u(gx)$ and $u(x)$ are uniformly close whenever both points are outside
$O$. Write $P_\epsilon1_O(x)=\int_G1_O(gx)\rho_\epsilon(g)\,\mathrm dg$.
For small kernel support, outside $O$ the set
$\{|u_\epsilon-u|>\delta\}$ is contained in
$\{P_\epsilon1_O>\delta/(4L)\}$. Markov and invariance give
\[
 \operatorname{Cap}_\omega
    \{P_\epsilon1_O>\delta/(4L)\}
 \leq\frac{4L}{\delta}
       \sup_v\int_XP_\epsilon1_O\,\MA_\omega(v)
 \leq\frac{4L}{\delta}\operatorname{Cap}_\omega(O).
\]
Together with $\operatorname{Cap}_\omega(O)$ this proves capacity
convergence. If $L=0$, the assertion is immediate.

Here is the product passage for any uniformly bounded sequence $u_j$.
In a corrected relatively compact chart put $p_j=u_j+\rho$ and
$p=u+\rho$, genuine uniformly bounded local psh functions. The capacity
comparisons give local capacity convergence. For every compactly supported
smooth test form, telescoping the closed products and bounded integration
by parts express their difference as a sum of integrals of
\[
 (p_j-p)\ddc(\text{test})\wedge
          (\ddc p_j)^l\wedge(\ddc p)^{k-1-l}.
\]
Their absolute values are bounded by $\int|p_j-p|\,\mathrm d\lambda_j$,
where positive mixed measures $\lambda_j$ include fixed Euclidean
quadratic factors. They have uniformly bounded total masses and
$\lambda_j(A)\leq C\operatorname{Cap}_{\rm loc}(A)$ on the support.
Indeed each mixed product occurs with a positive coefficient in the
top power of the sum of its bounded psh potentials and the quadratics;
normalize that sum to $[0,1]$. Splitting at $|p_j-p|\leq\delta$ proves
convergence in every mixed degree.

Write $\omega+\ddc u_j=\ddc p_j-\gamma$ with the fixed smooth
semipositive $\gamma=\ddc\rho-\omega$. Expanding the cube proves
bounded nonclosed product convergence. The form $\gamma$ need not be
closed; its powers are smooth test coefficients in the closed-product
argument. At bounded depth these positive top measures obey the original
bound $(1+L)^3\operatorname{Cap}_\omega(A)$.
A continuous depth weight is uniformly continuous on $[0,L]$, so its
varying error vanishes by capacity domination. For a fixed
quasi-continuous weight, extend its restriction outside an open set of
small capacity to a bounded continuous function; domination bounds the
error. Weak convergence against this continuous extension completes the
weighted product passage.
\end{proof}

The total estimate in \eqref{eq:smooth-upper} now holds for every bounded
invariant $u\leq0$. For its tail estimate take continuous depth cutoffs
majorizing $1_{\{-u>T\}}$, supported in $\{-u>T-\delta\}$, apply the
bounded weighted passage and the smooth estimate, then let
$\delta\downarrow0$. Continuity of $(1+\sqrt{2t})W(t)$ gives
$Q_{T-\delta}(W)\to Q_T(W)$.

\section{Unbounded potentials, uniform tails and continuity}
For every $u\in\mathcal R$ the canonical $u^K=\max(u,-K)$ is bounded,
genuine, invariant and has supremum zero. On $\{u>-K\}$ locality of
bounded products identifies its measure with the compatible finite
level defining $\mu_u$. The restricted measures increase to $\mu_u$.
Their weights agree with $W(-u)$ on these sets. Monotone convergence
therefore extends both estimates in \eqref{eq:smooth-upper} to every
unbounded $u\in\mathcal R$, with the same constants. This proves the
sufficient total and tail assertions.

For every Borel set $B$, splitting at depth $T$ and using the original
restricted-depth capacity estimate gives
\begin{equation}\label{eq:uniform-ac}
 \nu_{W,u}(B)\leq W(0)(1+T)^3\operatorname{Cap}_\omega(B)
                                      +12Q_T(W).
\end{equation}
If $\sqrt t\,W(t)\to0$, first take $T$ large and then capacity small;
this proves uniform absolute continuity.

Suppose $u_j\to u$ in capacity within $\mathcal R$. Fixed bounded
truncations $u_j^K$ converge in capacity to $u^K$. Their weighted
measures converge weakly by Lemma~\ref{lem:capacity}. To remove
truncation, outside $\{u\leq-K\}$ finite-level locality identifies
$\nu_{W,u^K}$ and $\nu_{W,u}$, while
\[
 \nu_{W,u^K}(\{u\leq-K\})\leq12Q_{K-1}(W),\qquad
 \nu_{W,u}(\{u\leq-K\})\leq12Q_{K-1}(W).
\]
The inequalities use the tail sets of $u^K$ and $u$, respectively,
and include possible bounded-truncation mass on $\{u=-K\}$.
Thus the total variation error is at most $24Q_{K-1}(W)$, uniformly
in all potentials. Letting $K\to\infty$ proves weighted weak
continuity when $\sqrt t\,W(t)\to0$.

\section{Capacity concentration and the critical discontinuity}
The core potentials depend only on $x,w$, and the decreasing sum used
for Theorem~\ref{thm:individual} preserves this invariance. Consequently
the single infinite measure and the necessary mass and tail conditions
already hold within $\mathcal R$.

The cores have $x\geq c_1/\eta$ for a fixed $c_1>0$, hence lie in
shrinking neighborhoods of the analytic hyperplane $H=\{z_0=0\}$.
It has zero original capacity by Lemma~\ref{lem:capacity}, and closed
neighborhoods of it have capacity tending to zero. If
$\limsup\sqrt t\,W(t)>0$, select $\eta$ so that
$D\eta^{-1}W(B\eta^{-2})$ is bounded below. The actual weighted
core masses violate uniform capacity absolute continuity. This also
covers unbounded $\sqrt t\,W(t)$.

For the failure of weak continuity in the bounded critical regime,
use the \emph{controlled} correction chosen in the individual-measure
proof. Its plateau and its late supplement satisfy
\[
 \int r^0_\eta=K\eta x_{\rm pre}+O(\eta),\qquad
 \int s_\eta=O(\sqrt{\eta L})+O(\eta),\qquad
 \eta x_{\rm pre}\to\ell.
\]
Therefore $m_\eta\to m:=K\ell>0$. Shrink the fixed $c$ so that $m$
satisfies all the established positivity margins. On each fixed finite
$x$ interval, $r_\eta$ and all its derivatives tend to zero; the late
supplement escapes the interval. Thus $g_\eta$ converges smoothly there to
\[
 g(x)=-m\int_{-\infty}^x(x-s)\psi(s)\,\mathrm ds.
\]
It is zero for $x\leq-2$ and equals $-m(x-\mu_\psi)$ for $x\geq-1$.
The same margins give $\mu+g'>0$ and
$g''+\mu(1-\mu)\geq0$ at finite $x$.
Its base extension at infinity is a canonical logarithmic pole $mH$.
It is therefore an actual normalized genuine $\omega$-psh potential
in $\mathcal R$.

The uniform heat bound \eqref{eq:offpolevalue} and the upper bound
$F_\eta\leq C\eta$ imply that, on compact base sets away from $H$,
fixed-depth exceptional sets lie in exponentially shrinking tubes about
the local holomorphic zero section $w=0$. Finite coordinate covers,
polar-set capacity zero and Lemma~\ref{lem:capacity} give capacity
tending to zero for these tubes. The corrections converge uniformly on
these base compact sets, and their supremum normalization is $O(\eta)$.
Remove an arbitrarily small capacity neighborhood of $H$. These
statements prove the actual convergence $u_\eta\to g$ in original
capacity, rather than convergence of formal energy profiles.

Assume $\sqrt t\,W(t)$ is bounded but does not tend to zero. Choose
the preceding sequence of cores with weighted mass bounded below by
$c_0>0$. The proved upper bound makes their full weighted measures
uniformly finite, so a subsequence has a weak cluster measure $\nu$.
For each fixed closed neighborhood $H_r$ of $H$, all sufficiently late
cores lie in $H_r$. Portmanteau gives $\nu(H_r)\geq c_0$; decreasing
continuity of the finite measure gives $\nu(H)\geq c_0$.
But $\nu_{W,g}(H)=0$ by its defining non-pluripolar truncation.
Thus weighted weak continuity fails for this actual capacity-convergent
sequence. If instead $\sqrt t\,W(t)$ is unbounded,
Theorem~\ref{thm:individual} already prevents a finite-measure map on
all of $\mathcal R$. These necessity arguments, the flux estimate and
the truncation passage complete Theorem~\ref{thm:sharp}.

\begin{remark}
The theorem treats every genuine potential in the stated symmetry class,
including unbounded ones; the counterexamples also obstruct the full
unrestricted class. Angular averaging alone supplies no upper bound for
the quadratic fiber energy of an arbitrary potential. No unrestricted
reduction or full solution of original S02 follows. All new conclusions
in this working article await independent mathematical review and
reconstruction. No literature-priority assertion is made.
\end{remark}

\begin{thebibliography}{9}
\bibitem{XZ}
M.~Xia and K.~Zhang,
\emph{A counterexample to the bounded mass property},
arXiv:2608.21053v1 (2026),
\url{https://arxiv.org/abs/2608.21053v1}.
\bibitem{BT}
E.~Bedford and B.~A.~Taylor,
\emph{A new capacity for plurisubharmonic functions},
Acta Math. 149 (1982), 1--40.
\end{thebibliography}
\end{document}
